\documentclass[UTF8,zihao=-4]{ctexart}

\usepackage[a4paper,margin=2.3cm]{geometry}
\usepackage{array}
\usepackage{booktabs}
\usepackage{longtable}
\usepackage{hyperref}
\usepackage{xcolor}
\usepackage{enumitem}
\usepackage{fancyhdr}
\usepackage{lastpage}

\hypersetup{
  colorlinks=true,
  linkcolor=black,
  urlcolor=blue
}

\setlength{\parindent}{2em}
\setlength{\parskip}{0.35em}
\setlength{\emergencystretch}{3em}
\setlength{\tabcolsep}{4pt}
\renewcommand{\arraystretch}{1.25}
\setlist[itemize]{leftmargin=2.2em,itemsep=0.2em,topsep=0.2em}
\pagestyle{fancy}
\fancyhf{}
\fancyfoot[C]{第 \thepage 页 / 共 \pageref*{LastPage} 页}
\renewcommand{\headrulewidth}{0pt}
\renewcommand{\footrulewidth}{0pt}

\title{\heiti 托福阅读 Class 6 课堂笔记}
\author{}
\date{}

\begin{document}
\maketitle
\thispagestyle{fancy}

本节课继续整理托福阅读中的同近义词群、词根词缀和题型功能词。内容包括“祖先、后代、描绘、雕刻、前沿”等高频词群，拉丁语和希腊语来源的不规则复数，\texttt{-scribe- / -script-}、\texttt{cosm-}、\texttt{heli-}、\texttt{-trude- / -trus-} 等核心词根，以及修辞目的题中常见的举例、支持、解释、反驳、强调、比较和定义类动词。

\section*{一、一般生词、同近义词群与高频短语}

表示“祖先、先驱、前身”的词群包括 \texttt{ancestor}、\texttt{predecessor}、\texttt{forerunner}、\texttt{precursor} 与 \texttt{antecedent}。其中 \texttt{predecessor} 可拆为 \texttt{pre-} “前” + 与“走、离去”相关的词根 + \texttt{-or} “人”，表示“走在前面的人”；\texttt{precursor} 可拆为 \texttt{pre-} “前” + \texttt{curs-} “跑” + \texttt{-or}，表示“跑在前面的人”；\texttt{antecedent} 中的 \texttt{ante-} 与 \texttt{ced-} 分别表示“前”和“走”。对应地，表示“子孙、后代”的词群包括 \texttt{progeny}、\texttt{posterity}、\texttt{offspring} 和 \texttt{descendant}。\texttt{progeny} 中的 \texttt{gen-} 与“生育、产生”有关，\texttt{posterity} 中的 \texttt{post-} 表示“后”，\texttt{descendant} 可理解为“沿血脉向下延续的人”。

描述与刻写类词汇可以放在一起记忆。\texttt{describe}、\texttt{depict} 和 \texttt{delineate} 都可表示“描述、描绘”，其中 \texttt{delineate} 可拆为 \texttt{de-} “完全” + \texttt{line} “线” + \texttt{-ate}，即“用线条勾勒出来”。\texttt{inscribe}、\texttt{carve}、\texttt{engrave} 和 \texttt{incise} 都与“雕刻、铭刻”有关，其中 \texttt{incise} 中的 \texttt{cis-} 与“切”有关，可联想到 \texttt{scissors}。

表示“前沿、先进”的表达包括 \texttt{at the cutting edge of}、\texttt{cutting-edge}、\texttt{at the forefront of} 和 \texttt{at the vanguard of}。\texttt{sophisticated} 在托福阅读中可表示“先进的、复杂的、精密的”，需要结合语境判断是技术层面的“先进精密”，还是思维或系统层面的“复杂成熟”。

其他高频词和搭配中，\texttt{give prominence to} 表示“突出……”，\texttt{rise to prominence} 表示“声名鹊起、崭露头角”；\texttt{hide} 作动词是“隐藏”，作名词可表示“兽皮”，而 \texttt{hideous} 表示“可怕的、极丑的”。\texttt{proliferation} 表示“激增、繁殖”，相关词有 \texttt{reproduction} “繁殖”、\texttt{prolific} “多产的”和 \texttt{productive} “高产的”。\texttt{constituent} 可作名词表示“成分、构成要素”，也可作形容词表示“构成的”；\texttt{pacify} 表示“安抚、平息”，其中 \texttt{pac-} 与“和平”有关，衍生词 \texttt{pacifier} 可表示“安抚奶嘴、平息者”。\texttt{deferential} 表示“恭敬的、顺从的”，\texttt{uninhabited} 表示“无人居住的”，\texttt{a constellation of} 表示“一系列、一群”，\texttt{in proportion to} 和 \texttt{in inverse proportion to} 分别表示“与……成比例”和“与……成反比”。

\section*{二、托福高频不规则复数}

托福阅读中经常出现拉丁语、希腊语来源的不规则复数形式。整理时需要把单数、复数和语义一起记住，尤其是自然科学文章中的 \texttt{nucleus / nuclei}、\texttt{stimulus / stimuli}、\texttt{fungus / fungi}、\texttt{alga / algae}、\texttt{larva / larvae} 等词。

\begin{longtable}{>{\raggedright\arraybackslash}p{3.2cm} >{\raggedright\arraybackslash}p{4.2cm} >{\raggedright\arraybackslash}p{4.2cm} >{\raggedright\arraybackslash}p{3.3cm}}
\toprule
\textbf{规则} & \textbf{单数} & \textbf{复数} & \textbf{含义} \\
\midrule
\endfirsthead
\toprule
\textbf{规则} & \textbf{单数} & \textbf{复数} & \textbf{含义} \\
\midrule
\endhead
\texttt{-us -> -i} & \texttt{nucleus} & \texttt{nuclei} & 细胞核；核心 \\
 & \texttt{stimulus} & \texttt{stimuli} & 刺激物 \\
 & \texttt{alumnus} & \texttt{alumni} & 校友 \\
 & \texttt{fungus} & \texttt{fungi} & 真菌 \\
 & \texttt{cactus} & \texttt{cacti} & 仙人掌 \\
 & \texttt{locus} & \texttt{loci} & 地点；轨迹 \\
 & \texttt{syllabus} & \texttt{syllabi} & 教学大纲 \\
\texttt{-on / -um -> -a} & \texttt{protozoon} & \texttt{protozoa} & 原生动物 \\
 & \texttt{spectrum} & \texttt{spectra} & 光谱；范围 \\
 & \texttt{quantum} & \texttt{quanta} & 量子 \\
 & \texttt{millennium} & \texttt{millennia} & 千禧年；一千年 \\
\texttt{-a -> -ae} & \texttt{alga} & \texttt{algae} & 藻类 \\
 & \texttt{formula} & \texttt{formulae} & 公式；配方 \\
 & \texttt{larva} & \texttt{larvae} & 幼虫 \\
 & \texttt{nebula} & \texttt{nebulae} & 星云 \\
\texttt{-is -> -es} & \texttt{crisis} & \texttt{crises} & 危机 \\
\bottomrule
\end{longtable}

\section*{三、核心词根：写、世界与太阳}

\texttt{-scribe- / -script-} 表示“写、画”。\texttt{describe} 是“描述”，可理解为“写下来”；\texttt{inscribe} 表示“雕刻、铭刻”，相关名词为 \texttt{inscription}，搭配 \texttt{be inscribed with meaning} 可表示“被赋予意义”；\texttt{prescribe} 表示“开处方、规定”，名词为 \texttt{prescription}，形容词 \texttt{prescriptive} 表示“规定性的”。\texttt{subscribe to} 既可表示“订阅”，也可表示“赞同”，词源意象是“在文件下方签名表示同意”；\texttt{ascribe sth to sth.} 表示“把……归因于……”；\texttt{transcribe} 表示“转录、抄写”，\texttt{transcript} 可表示“成绩单、文字记录”；\texttt{manuscript} 表示“手稿”，其中 \texttt{manu-} 与“手”有关；\texttt{scribble} 表示“涂鸦、潦草地写”，\texttt{script} 表示“剧本、脚本”。

\texttt{cosm-} 表示“宇宙、世界”。\texttt{cosmos} 表示“宇宙”，\texttt{cosmic} 表示“宇宙的”；\texttt{cosmopolitan} 可作形容词表示“世界性的”，也可作名词表示“世界主义者”，其中 \texttt{polis} 与“城市”有关；\texttt{metropolitan} 表示“大都市的”，可联系 \texttt{metro-} 与 \texttt{polis} 记忆。

\texttt{heli- / helio-} 表示“太阳”。\texttt{helium} 是“氦气”，因为其最初是在太阳光谱中被发现；\texttt{heliocentric} 表示“日心的”，如日心说；\texttt{heliotherapy} 表示“日光浴疗法”；\texttt{heliosis} 表示“中暑”，其中 \texttt{-osis} 常作疾病后缀；\texttt{heliotropic} 表示“向日性的”，其中 \texttt{trop-} 表示“转向”。

\section*{四、托福阅读修辞目的题功能词}

修辞目的题常问 \texttt{Why does the author mention/discuss/refer to ... ?}，选项中经常出现表示“举例、支持、解释、反驳、强调、对比、定义”的功能动词。掌握这些词有助于快速判断句子或段落在文章中的作用。

表示“举例与展示”的词包括 \texttt{exemplify}、\texttt{illustrate} 和 \texttt{demonstrate}。\texttt{exemplify} 表示“举例说明、作为典型”，\texttt{illustrate} 本义与“照亮、加插图”有关，引申为“说明”，\texttt{demonstrate} 表示“证明、展示”。表示“支持与强化”的表达包括 \texttt{support}、\texttt{back up the claim that}、\texttt{provide evidence that} 和 \texttt{reinforce the point that}，其中 \texttt{reinforce} 可理解为“再次注入力量”，即“强化、加固”。

表示“解释与阐明”的词包括 \texttt{account for}、\texttt{explain}、\texttt{elaborate on}、\texttt{shed light on} 和 \texttt{clarify}。\texttt{account for} 既可表示“解释原因”，也可表示“占……比例”；\texttt{elaborate on} 表示“详细阐述”；\texttt{shed light on} 字面义是“把光照在……上”，引申为“阐明、提供线索”；\texttt{clarify} 表示“澄清、阐明”。表示“反驳与质疑”的词包括 \texttt{refute}、\texttt{contradict}、\texttt{criticize}、\texttt{question}、\texttt{challenge}、\texttt{argue against} 和 \texttt{cast doubt on}，其中 \texttt{contradict} 可拆为 \texttt{contra-} “相反” + \texttt{dict-} “说”，即“说出相反的话”。

表示“强调与引起注意”的词包括 \texttt{highlight}、\texttt{stress}、\texttt{emphasize}、\texttt{underscore}、\texttt{draw attention to}、\texttt{call attention to} 和 \texttt{point out}。\texttt{underscore} 原指“在下面划线”，后来引申为“强调”。表示“比较与对比”的词包括 \texttt{contrast}、\texttt{compare} 和 \texttt{distinguish between ...}；表示“定义与特征”的词包括 \texttt{define}、\texttt{identify}、\texttt{characterize} 和 \texttt{specify}，其中 \texttt{specify} 中的 \texttt{spec-} 与“看、具体”有关，表示“具体说明、明确指出”。

\section*{五、气候地理词、侵入词群与 \texttt{-trude-} 词根}

气候与地理类词汇中，\texttt{trade wind} 表示“信风、贸易风”，常见于地理和历史类文章；\texttt{precipitation} 表示“降水”，包括雨、雪、冰雹等，也可表示“沉淀物”；\texttt{mass transportation} 或 \texttt{mass transit} 表示“公共交通、大运量交通”，\texttt{mass media} 表示“大众媒体”。\texttt{solicit} 表示“征求、恳求、筹集”，可用于 \texttt{solicit donations}；\texttt{vibrant} 表示“充满活力的、色彩鲜艳的、声音洪亮的”，其中 \texttt{vibr-} 与“振动”有关。

表示“侵入、侵占”的词群包括 \texttt{invade}、\texttt{invasion}、\texttt{incursion} 和 \texttt{encroach on}。\texttt{invade} 强调“入侵、侵略”，常用于武装力量或大规模涌入；\texttt{invasion} 是其名词形式。\texttt{incursion} 表示“突然的入侵”，其中 \texttt{curs-} 与“跑”有关，可理解为“跑进别人的地盘”；\texttt{encroach on} 表示“蚕食、侵占”，常指土地、权利、时间等被逐渐侵入。

\texttt{-trude- / -trus-} 表示“推”。\texttt{intrude} 表示“侵入、闯入、打扰”，可拆为 \texttt{in-} “进入” + \texttt{-trude-} “推”，常见搭配为 \texttt{intrude on / upon}；\texttt{intrusion} 是“侵入、打扰”。\texttt{protrude} 表示“突出、伸出、鼓出”，其中 \texttt{pro-} 表示“向前”；\texttt{extrude} 表示“挤出、压出”，在制造业或地质学中可指物质从内部被挤压到外部；\texttt{obtrude} 表示“强加于人、使突出”，常带贬义，搭配 \texttt{obtrude oneself} 可表示“强出头、到处插手”。

\section*{六、补充词汇与固定表达}

\begin{longtable}{>{\raggedright\arraybackslash}p{5.0cm} >{\raggedright\arraybackslash}p{9.5cm}}
\toprule
\textbf{单词或表达} & \textbf{含义或用法} \\
\midrule
\endfirsthead
\toprule
\textbf{单词或表达} & \textbf{含义或用法} \\
\midrule
\endhead
\texttt{ancestor / predecessor / precursor} & 祖先；前辈；先驱；前身 \\
\texttt{progeny / posterity / offspring} & 子孙，后代 \\
\texttt{describe / depict / delineate} & 描述，描绘，勾勒 \\
\texttt{inscribe / carve / engrave / incise} & 雕刻，铭刻，切入 \\
\texttt{cutting-edge} & 尖端的，前沿的 \\
\texttt{at the forefront of} & 在……的最前沿 \\
\texttt{give prominence to} & 突出…… \\
\texttt{rise to prominence} & 崭露头角，声名鹊起 \\
\texttt{proliferation} & 激增；繁殖 \\
\texttt{constituent} & 成分，构成要素；构成的 \\
\texttt{pacify / pacifier} & 安抚、平息 / 安抚奶嘴；平息者 \\
\texttt{deferential} & 恭敬的，顺从的 \\
\texttt{uninhabited} & 无人居住的 \\
\texttt{a constellation of} & 一系列；一群 \\
\texttt{in inverse proportion to} & 与……成反比 \\
\texttt{inscription / prescription} & 铭文 / 处方；规定 \\
\texttt{subscribe to / ascribe to} & 订阅、赞同 / 归因于 \\
\texttt{transcribe / transcript} & 转录、抄写 / 成绩单；文字记录 \\
\texttt{cosmos / cosmic} & 宇宙 / 宇宙的 \\
\texttt{cosmopolitan / metropolitan} & 世界性的；世界主义者 / 大都市的 \\
\texttt{heliocentric / heliotropic} & 日心的 / 向日性的 \\
\texttt{exemplify / illustrate} & 举例说明 / 说明、阐明 \\
\texttt{reinforce the point that} & 强化某一观点 \\
\texttt{elaborate on / clarify} & 详细阐述 / 澄清、阐明 \\
\texttt{refute / contradict} & 反驳 / 与……矛盾 \\
\texttt{cast doubt on} & 对……产生怀疑 \\
\texttt{underscore / highlight} & 强调；使显著 \\
\texttt{distinguish between ...} & 区分…… \\
\texttt{characterize / specify} & 描绘特征 / 具体说明 \\
\texttt{trade wind / precipitation} & 信风 / 降水；沉淀物 \\
\texttt{mass transit / mass media} & 公共交通 / 大众媒体 \\
\texttt{solicit / vibrant} & 征求、筹集 / 充满活力的 \\
\texttt{invade / incursion / encroach on} & 入侵 / 突然侵入 / 蚕食、侵占 \\
\texttt{intrude / intrusion} & 闯入、打扰 / 侵入、打扰 \\
\texttt{protrude / extrude / obtrude} & 突出、伸出 / 挤出、压出 / 强加于人 \\
\bottomrule
\end{longtable}

\end{document}
